\documentclass[11pt,a4paper]{article}

\usepackage{kotex}
\usepackage{fontspec}
\setmainfont{Times New Roman}
\setmainhangulfont{AppleMyungjo}
\setsanshangulfont{Apple SD Gothic Neo}

\usepackage[a4paper,top=18mm,bottom=18mm,left=20mm,right=20mm]{geometry}
\usepackage{fancyhdr}
\setlength{\headheight}{24pt}
\usepackage{enumitem}
\usepackage{mdframed}
\usepackage{array}
\usepackage{booktabs}

\setlength{\parindent}{1em}
\linespread{1.18}

\pagestyle{fancy}
\fancyhf{}
\renewcommand{\headrulewidth}{0.6pt}
\fancyhead[L]{\sffamily\bfseries 2026 고1 영어 변형 --- 지문 36}
\fancyhead[R]{\sffamily 빙하기와 농업의 시작}
\renewcommand{\footrulewidth}{0pt}
\fancyfoot[C]{\framebox[16mm][c]{\thepage}}

\newcommand{\qno}[1]{\par\vspace{8pt}\noindent{\sffamily\bfseries #1.}\ }
\newcommand{\instr}[1]{{\sffamily #1}\par\vspace{3pt}}
\newlist{choices}{enumerate}{1}
\setlist[choices]{label=,leftmargin=1.5em,itemsep=2pt,topsep=3pt,parsep=0pt}
\newmdenv[linewidth=0.6pt,roundcorner=3pt,innertopmargin=7pt,innerbottommargin=7pt,
          innerleftmargin=9pt,innerrightmargin=9pt,skipabove=5pt,skipbelow=5pt]{pbox}
\newcommand{\nemo}[1]{\fbox{\,#1\,}}

\begin{document}

\begin{center}
{\fontsize{15}{17}\selectfont\sffamily\bfseries [지문 36] 다음 글을 읽고 물음에 답하시오.}
\end{center}
\vspace{4pt}

% ===== 주어진 단락 (Opening paragraph in pbox) =====
\begin{pbox}
During the Ice Age, not only was the planet much colder, with ice sheets covering much of what we call the northern hemisphere today, but, crucially, it was much drier.
\end{pbox}

\vspace{6pt}

% ===== (A)(B)(C) 섹션 (pbox 바깥, 들여쓰기 + 볼드 라벨) =====
\noindent\textbf{(A)}\quad In this type of climate, farming isn't an option: it's too risky to depend on any one piece of land to produce the energy you need. As the temperature rose and the ice caps melted, we experienced a sudden explosion of life.

\vspace{6pt}

\noindent\textbf{(B)}\quad In Ireland, we often associate the cold with the wet but if it is really cold, there is far less evaporation$^{*}$, fewer clouds, and less rain. Our world in the Ice Age was cold and dry, meaning it was difficult for plants to grow.

\vspace{6pt}

\noindent\textbf{(C)}\quad The world got warmer and wetter, and people started to live around places where they could make food grow most intensively. This didn't happen overnight; it probably took thousands of years, with hunter-gatherers foraging$^{**}$ and hunting, while doing a little side hustle$^{***}$ in farming.

\vspace{4pt}
\noindent\footnotesize{$^{*}$ evaporation: 증발 \quad $^{**}$ forage: 채집하다 \quad $^{***}$ side hustle: 부업}
\normalsize

\bigskip

% ===================== Q1 =====================
\qno{1}\instr{주어진 글 다음에 이어질 글의 순서로 가장 적절한 것은?}
\begin{choices}
\item ① (A) -- (C) -- (B)
\item ② (B) -- (A) -- (C)
\item ③ (B) -- (C) -- (A)
\item ④ (C) -- (A) -- (B)
\item ⑤ (C) -- (B) -- (A)
\end{choices}

\bigskip

% ===================== Q2 =====================
\qno{2}\instr{(B) 단락의 논리 전개로 가장 적절한 것은?}
\begin{choices}
\item ① Very low temperatures weaken the movement of moisture into the air, reducing cloud formation and rainfall.
\item ② Low temperatures make Ireland unusually wet, so farming becomes possible despite poor soil.
\item ③ A lack of clouds raises temperatures, causing ice sheets to melt and rainfall to increase.
\item ④ Less rain is presented as the result of farming, not as a condition that limits plant growth.
\item ⑤ Cold weather is described as useful because it keeps plants from using too much energy.
\end{choices}

\bigskip

% ===================== Q3 (서술형) =====================
\qno{3}\instr{[서술형] (C) 단락에서 농업이 갑자기 시작된 것이 아님을 나타내는 영어 표현을 찾아 쓰시오.}

\vspace{6pt}
\noindent $\Rightarrow$ \rule{10cm}{0.4pt}

% ===================== Q4 =====================
\qno{4}\instr{윗글의 제목으로 가장 적절한 것은?}
\begin{choices}
\item ① From Cold Dry Lands to the Conditions for Farming
\item ② Why Ice Sheets Made Ireland Warmer and Wetter
\item ③ The Sudden End of Hunting in the Northern Hemisphere
\item ④ Farming as the Only Reliable Strategy During the Ice Age
\item ⑤ How Clouds Prevented Humans from Gathering Food
\end{choices}

\end{document}
